\documentclass[11pt]{article}

\usepackage[utf8]{inputenc}
\usepackage[T1]{fontenc}
\usepackage[margin=1in]{geometry}
\usepackage[strings]{underscore}
\usepackage{hyperref}
\usepackage{titlesec}
\usepackage{tikz}
\usetikzlibrary{calc}
\usepackage{eso-pic}

\AddToShipoutPictureBG{%
  \begin{tikzpicture}[remember picture,overlay]
    \draw[line width=0.8pt]
      ($(current page.north west)+(0.5in,-0.5in)$) rectangle
      ($(current page.south east)+(-0.5in,0.5in)$);
  \end{tikzpicture}%
}

\hypersetup{
  pdftitle={NOMAD: Node Out-of-Band Module for Attestation \& Defense -- Abstract},
  pdfauthor={Kernel Borderlands Team},
  colorlinks=true,
  linkcolor=black,
  citecolor=black,
  urlcolor=blue
}

\titleformat{\section}{\normalfont\bfseries}{\thesection}{0.5em}{}
\setlength{\parskip}{0.5em}
\setlength{\parindent}{0pt}
\renewcommand{\arraystretch}{1.05}

\title{\textbf{NOMAD: A Node Out-of-Band Module for Attestation and Defense}}

\author{Pardhu Sri Rushi Varma K \\ \small 2311CS040089}
\date{}

\begin{document}
\pagestyle{empty}

% \author is kept above for record-keeping but intentionally not rendered:
% title is typeset manually so the author line never appears in the PDF.
\makeatletter
\begin{center}
{\LARGE\bfseries \@title \par}
\end{center}
\makeatother
\vspace{-1em}

\begin{center}
\small Department of Cybersecurity, ECE \& IoT, Malla Reddy University, Hyderabad
\end{center}

\vspace{0.5em}

\section*{Abstract}

Kernel-level runtime defenses face an inherent limit: any detector that runs
as software on the host it protects can be blinded or disabled by an
attacker who gains sufficient privilege on that same host. Once kernel
telemetry, logging, and enforcement all depend on the honesty of the kernel
being monitored, a full kernel compromise removes the very evidence needed
to detect it. We present NOMAD (Node Out-of-Band Module for Attestation and
Defense), a hardware-based attestation and containment module for the
Kernel Borderlands runtime defense framework. It runs on a physically
separate, minimal computing complex attached to the host, and communicates
with the host's software-based watchdog over a dedicated out-of-band link
rather than through the monitored operating system. It periodically attests
to the host's integrity and, when attestation fails or a containment
decision is issued, can act through a hardware relay capable of cutting the
host's power independent of the host's own software state. This places the
last line of defense outside the failure domain it is meant to defend,
so that a full kernel compromise -- including rootkits with Ring-0
privilege -- does not also compromise the ability to detect and contain it.
We describe the design rationale, the threat model it addresses,
the two isolated communication planes it exposes, and its current status
as an architectural design within the broader Kernel Borderlands project.

\vspace{0.7em}
\noindent\textbf{Keywords:} kernel security, out-of-band attestation, root
of trust, hardware watchdog, runtime defense, rootkit resistance

\vspace{0.7em}
\noindent\textbf{Team Members:}

\small
\begin{center}
\begin{tabular}{@{}ll@{\hspace{2em}}ll@{}}
Md Vaseem & 2311CS040004 & Tejaswini K & 2311CS040077 \\
Harshitha B & 2311CS040022 & Vamshi K & 2311CS040078 \\
Sai Raghava D & 2311CS040041 & Karthik K & 2311CS040080 \\
Ashwin Kumar D & 2311CS040047 & Rupa Yeshvitha K & 2311CS040082 \\
Madhumitha J & 2311CS040067 & Rajesh K & 2311CS040084 \\
Jeshwanth K & 2311CS040072 & Pardhu Sri Rushi Varma K & 2311CS040089 \\
Agasthya K & 2311CS040075 & & \\
\end{tabular}
\end{center}

\end{document}
